\chapter{Future works}
\label{cap:future_works}

\section{Other functionalities and integrations}

Looking ahead, different directions for future developments have been individuated to enhance both functionality and user experience through the framework.

A possible feature to add is the integration of a Time Capsule system, enabling data to be stored in a small physical sensors distributed across urban environments, natural landscapes or remote places. Extending the application to integrate sensors such as Bluetooth or Near Field
Communication (NFC) which working as short-range communication technologies thus to bring users physically closer to specific locations. This functionality opens the app to a vast scenarios such as gameplay, for example an interactive escape room or an enhanced scavenger hunts, but also works as a memories holder, storing photographs, videos and significant content.

Altitude is a metadata currently stored during the content creation, but not yet implemented through the map. By using this dimension, GeoMedia, could support vertical placement and filtering, adaptable to high-rise buildings, mountains or drone-assisted frameworks.

Another functionality not currently implemented by GeoMedia is the integration of Augmented Reality operations. However, the app provide an in-app camera component that can be leveraged for future AR-based enhancements. This feature could enable the integration of digital information above the real-word environment, visualizing for example historical information or media content. With this development, the mobile application would improve the user experience by providing a more immersive and interactive way to explore and interact with geographically referenced content.

While these extensions primarily concern the functional capabilities of GeoMedia, future developments could also involve changes to the underlying system architecture. As mentioned previously, an improvement in this direction is the adoption of microservices-based infrastructure, and thus allowing the OSN to operate world-wide without suffering in terms of performance and responsivity; while promoting scalability, fault tolerance and maintainability and aligning with framework with modern cloud-native platforms.

\section{Mesh network}
\label{sub:nostr}
Furthermore, future versions of GeoMedia could implement a mesh network to support a more decentralized and potentially a serverless architecture; allowing users to share content based on their physical proximity.
By leveraging protocols such as Bluetooth Low Energy (BLE), or standard Bluetooth, the GeoMedia framework could establish peer-to-peer connections between nearby devices without requiring continuous access to the Internet or cellular infrastructure. A similar approach is implemented by Bitchat~\cite{site:bitchat}, a messaging application that enables users to exchange messages over Bluetooth without requiring Internet connectivity, cellular service, user accounts or centralized servers. \\
As Bitchat does, even GeoMedia, in addition could adopt the local peer-to-peer Nostr~\cite{site:nostr} framework, an open and decentralized protocol that allows user to exchange information and data through distributed relays. Nostr could complement the local mesh layer by providing a synchronizing mechanism when Internet connectivity is available. In example, users could create and publish a posts related to an event, nearby devices could see it directly over the mesh network, while others users could subsequently synchronize the post. This scenario create a challenge of coordinate an hybrid architecture in which BLE communication provide local content meanwhile Nostr provides decentralized synchronization to further distance up to the maximum range imposed by GeoMedia (currently \texttt{30KM}). As result, this solution would reduce reliance on a single central server, improving the resilience of GeoMedia instance in environment where Internet connectivity is intermitted, unavailable or intentionally restricted. 
However, implementing this technology would introduce several challenges related to data consistency, content storage, privacy, message propagation and resource resource consumption on mobile devices.

\section{Civil Applications and Social Impact}

GeoMedia also presents an opportunity to encourage and incentivize civic participation around local events and news. Platforms such as Waze~\cite{site:waze} focus primarily on location-based navigation and driving assistance, demonstrating how geographically contextualized information can be used to provide users with relevant content and services. GeoMedia could extend this concept beyond navigation by enabling grassroots content creation and dissemination, allowing communities to document events near them, express collective concerns or to simply share geographically contextualized knowledge.\\
In a democratic scenario, these capabilities may contribute to greater civic awareness, transparency and participatory governance. Instead, as opposite, in environments characterized by censorship or centralized control of information, features such as decentralized media storage and location-bound data access could potentially provide alternative channels for information dissemination and access. This aspect would make GeoMedia particularly relevant for investigating the role of location-based systems in supporting information resilience and civic communication under restrictive conditions. However, these capabilities also raise important concerns regarding privacy, security or fake-news and misinformation. This feature would bring GeoMedia to an interdisciplinary level, examining both the civic benefits than societal risks associated with the framework, with attention to topics such as the digital freedom, censorship and ethical usage of location-based systems.
